\documentclass[11pt]{article}
\usepackage[a4paper,margin=27mm]{geometry}
\usepackage{amsmath,amssymb,amsthm,hyperref}
\emergencystretch=2em
\newtheorem{theorem}{Theorem}
\newtheorem{lemma}{Lemma}
\newcommand{\E}{\mathbb E}
\newcommand{\Pol}{\operatorname{Pol}}
\title{A power lower bound for every permutation-fair die}
\author{Research proof: independently reviewed}
\date{30 September 2026}
\begin{document}\maketitle
\noindent\textbf{Status.} Complete proof with independent mathematical reviews by the three
other research agents. These are not formal proof-assistant checks.
Its essential input is the finite mean-deficit balance
theorem in
\path{rational_designs/asymptotic/finite_endpoint_odds_balance.tex}.

\begin{theorem}\label{main}
There is an absolute constant $c>0$ such that every individual die in
a permutation-fair collection of $n$ equiprobable finite dice has at
least $c n^{3/2}$ faces.
More generally, suppose $p_1,\ldots,p_m:[K]\to[0,1]$ are nondecreasing,
the rows have positive probability weights $w_q$, and
\[
 \E_w\prod_{j\in S}p_j=\frac1{|S|+1}\qquad(S\subseteq[m]).
\]
Then $w_1,w_K\le C m^{-3/2}$ for an absolute $C$.
\end{theorem}

All constants denoted by $C,c$ below are absolute and may change from
line to line. Fixed numerical choices, such as $\theta=1/16$, do not
depend on $m$ or on the representation. We prove the assertion for
sufficiently large $m$; increasing $C$ covers the remaining range.

\section*{The finite balance input}
Put $u_j=1-p_j$ and $s=\sum_j u_j$. The cited finite balance theorem
and its row-mass corollary give, for $m\ge256$,
\begin{align}
 \max_j u_j&\le256(s+1)/m,                    \label{balance}\\
 \Pr_w(s\le S)&\le512(S+1)/m\quad(S\ge0).          \label{mass}
\end{align}
The first bound is the cited theorem when $s\le m/256$; otherwise
it follows simply from $u_j\le1$. In particular, for every fixed $a>0$,
\begin{equation}\label{laplace}
 \E_w e^{-as}\le C_a/m.
\end{equation}
We shall also use the following uniform consequence for integers $k\ge0$:
\begin{equation}\label{weightedmoments}
 \E_w (s+1)^k e^{-as}\le \frac{C_a^{k+1}k!}{m}.
\end{equation}
Indeed split $e^{-as}=e^{-as/2}e^{-as/2}$, apply \eqref{laplace}, and
use $\sup_{s\ge0}(s+1)^k e^{-as/2}\le C_a^k k!$.

\section*{A stable polarization comparison}
For a polynomial $P(z)=\sum_k a_kz^k$ of degree at most $N$, write
\[
 \Pol_N(P)(z_1,\ldots,z_N)
    =\sum_k a_k\frac{e_k(z_1,\ldots,z_N)}{\binom Nk}.
\]
This is its symmetric multiaffine polarization.

\begin{lemma}[Centered elementary coefficients]\label{elementary}
If $\sum_{i=1}^Ny_i=0$ and $\sigma^2=\sum_i y_i^2$, then for
$1\le k\le N$,
\[
 \frac{|e_k(y)|}{\binom Nk}
       \le \left(\frac{\sqrt{e k}\,\sigma}{N}\right)^k.
\]
\end{lemma}
\begin{proof}
The case $\sigma=0$ is immediate. The elementary inequality
$|1+u|^2\le\exp(2\operatorname{Re}u+|u|^2)$ holds for every complex $u$.
Since the linear terms sum to zero, it gives
\[
 \left|\prod_i(1+zy_i)\right|\le\exp(\sigma^2|z|^2/2)
\]
for every complex $z$. Cauchy's coefficient bound at radius
$\sqrt{k}/\sigma$ gives $|e_k|\le(\sqrt e\,\sigma/\sqrt k)^k$.
Divide by $\binom Nk\ge(N/k)^k$.
\end{proof}

\begin{lemma}[A variance-adapted Laguerre derivative bound]\label{lagderiv}
For integers $j\ge h\ge0$ and $z>0$,
\[
 z^{h/2}\left|\frac{d^h}{dz^h}L_j(\theta z)\right|
       \le(\theta j)^{h/2}e^{\theta z/2}.
\]
For $h>j$ the derivative is zero.
\end{lemma}
\begin{proof}
Give entire functions the Gaussian norm
$\|f\|^2=\pi^{-1}\int_{\mathbb C}|f(w)|^2e^{-|w|^2}\,dA(w)$.
The functions $f_j(w)=w^j/\sqrt{j!}$ are orthonormal. The displacement
\[
 (D_af)(w)=e^{-|a|^2/2+aw}f(w-\overline a)
\]
is an isometry, by the change of variables $w=u+\overline a$.
For real $a=\sqrt y$, direct coefficient expansion gives
\[
 \left|\langle f_{j-h},D_af_j\rangle\right|
   =e^{-y/2}\sqrt{\frac{(j-h)!}{j!}}\,
            y^{h/2}\left|L_{j-h}^{(h)}(y)\right|\le1.
\]
Here the superscript in $L_{j-h}^{(h)}$ is its associated parameter;
indeed the coefficient of $w^{j-h}$ in $e^{aw}(w-a)^j$ is
$(-a)^h L_{j-h}^{(h)}(a^2)$.
The differentiation identity
$\frac{d^h}{dy^h}L_j(y)=(-1)^hL_{j-h}^{(h)}(y)$ and
$j!/(j-h)!\le j^h$ prove the unscaled assertion.
Apply the chain rule with $y=\theta z$.
\end{proof}

Fix from now on
\[
 \ell=\lfloor m/4\rfloor,\quad N=m-\ell,\quad
 \alpha=\ell/m,\quad \theta=1/16.
\]
Let $B$ be a uniformly chosen $\ell$-element subset of $[m]$, and put
\[
 A_B=\prod_{j\in B}p_j,\qquad \overline A=\E_B A_B,
 \qquad x=\alpha s.
\]
The measure $\mu$ on rows defined by
\[
 \mu(q)=(\ell+1)w_q\overline A(q)
\]
is a probability measure. Define also the probability density
\[
 b_\ell(z)=\frac{\ell+1}{\ell}(1-z/\ell)^\ell
                 \boldsymbol1_{[0,\ell]}(z).
\]
Its expectation functional is denoted by $\beta_\ell$.

Let $L_j$ be the ordinary Laguerre polynomials, orthonormal for
$e^{-z}\,dz$ on $[0,\infty)$.

\begin{lemma}[Comparison with a scalar endpoint measure]\label{comparison}
Let $1\le d\le\sqrt m$ and
\[
 Q_a(z)=\sum_{j=0}^d c_{aj}L_j(\theta z),\qquad
 B_a=\sum_j|c_{aj}|\quad(a=1,2).
\]
For $P=Q_1Q_2$, one has
\[
 \left|\E_\mu P(x)-\beta_\ell(P)\right|
       \le C B_1B_2\left(\frac d m+e^{-cm}\right).
 \tag{*}
\]
The constants are uniform in the coefficients.
\end{lemma}
\begin{proof}
Lemma~\ref{lagderiv} and the product rule give, for $z>0$,
\begin{equation}\label{derivatives}
 \frac{|P^{(k)}(z)|}{k!}
 \le B_1B_2 e^{\theta z}\frac{(2\sqrt{\theta d/z})^k}{k!}.
\end{equation}
We will occasionally also use the coarser, nonsingular derivative
bounds $|P'(z)|\le C B_1B_2d e^{\theta z}$ and
$|P(z)|\le B_1B_2e^{\theta z}$. The former follows from
$|L_j'(y)|\le j e^{y/2}$, using
$L_j'=-\sum_{v<j}L_v$ and the $h=0$ case of Lemma~\ref{lagderiv}.

For fixed $B$, put $z_j=\ell u_j$ for $j\notin B$, and let
$v_B=N^{-1}\sum_{j\notin B}z_j$ be their mean.
Multilinearity and the exact mixed moments give
\begin{equation}\label{exactpolarization}
 (\ell+1)\E_w\E_B\left[A_B\Pol_N(P)(z_{[m]\setminus B})\right]
       =\beta_\ell(P).
\end{equation}
Indeed, after expansion every distinct-coordinate product can be
replaced by the corresponding product of one common uniform variable
$t$. The right side becomes
$(\ell+1)\int_0^1t^\ell P(\ell(1-t))\,dt$.
The degree condition holds because $\deg P\le2d\le N$ for large $m$.

On every row, \eqref{balance} gives $\max_jz_j\le C(s+1)$, while
\[
 0\le x,v_B\le\frac{\alpha}{1-\alpha}s,
 \qquad
 \sigma_B^2:=\sum_{j\notin B}(z_j-v_B)^2
       \le\sum_{j\notin B}z_j^2\le C N(s+1)v_B.
\]
Maclaurin's inequality gives
\begin{equation}\label{abar}
 \overline A=\frac{e_\ell(p)}{\binom m\ell}
       \le(1-s/m)^\ell\le e^{-\alpha s}.
\end{equation}
If $v_B=0$, all its nonnegative $z_j$ are zero and polarization has
no error. Otherwise Taylor expansion about $v_B$ is exact:
\[
 \Pol_N(P)(z)-P(v_B)
 =\sum_{k=2}^{\deg P}\frac{P^{(k)}(v_B)}{k!}
                 \frac{e_k(z-v_B)}{\binom Nk}.
\]
The factors $v_B^{-k/2}$ in \eqref{derivatives} cancel the factors
$v_B^{k/2}$ in Lemma~\ref{elementary}. Thus the absolute difference
is at most
\[
 B_1B_2e^{\theta v_B}
       \sum_{k=2}^{2d}\frac{(C\sqrt{dk(s+1)/m})^k}{k!}.
\]
As $\theta\alpha/(1-\alpha)<\alpha/4$,
$\E_B A_B e^{\theta v_B}\le e^{-3\alpha s/4}$.
Using \eqref{weightedmoments}, with $\lceil k/2\rceil$ in place of
$k$, and multiplying by $\ell+1$, bounds this part of the integrated
error by
\[
 C B_1B_2\sum_{k=2}^{2d}
       (C\sqrt{d/m})^k
       \frac{k^{k/2}\lceil k/2\rceil!}{k!}.
\]
The last ratio is at most $C^k(k+1)$, as follows directly from
$k!\ge(k/e)^k$ and $\lceil k/2\rceil!\le\lceil k/2\rceil^{\lceil k/2\rceil}$.
For large $m$ and $d\le\sqrt m$, the resulting geometric series
is at most $Cd/m$. The polarization error is therefore
$CB_1B_2d/m$.

It remains to replace $P(v_B)$ by $P(x)$. Restrict first to
$s\le m/256$, where all $p_j\ge1/2$. Under the success-weighted law
$\Pr(B)=A_B/e_\ell(p)$, put $\Delta=v_B-x$.
The exact bias identity is
\[
 \E\Delta=\frac{\ell}{mN}\sum_{i<j}(p_i-p_j)^2
                   \frac{e_{\ell-1}(p_{-[i,j]})}{e_\ell(p)}.
\]
The ratio is at most $1/\min p_i\le2$. Moreover the subset-inclusion
covariances are nonpositive by Newton's inequality, giving
$\operatorname{Var}\Delta\le N^{-2}\sum_j(\ell u_j)^2$.
These facts are proved in detail in
\path{rational_designs/asymptotic/weighted_success_subset_moments.tex}.
Since
\[
 \sum_j u_j^2\le s\max_j u_j\le C s(s+1)/m,
\]
they give the sharper bounds
\begin{equation}\label{biasvariance}
 0\le\E\Delta\le C s(s+1)/m,\qquad
 \E\Delta^2\le C\bigl(s(s+1)/m+s^2(s+1)^2/m^2\bigr).
\end{equation}
If $s=0$, all the coordinates vanish and there is no averaging error.
Otherwise $x=\alpha s>0$. Taylor's exact integral remainder and
\eqref{derivatives} give
\[
 |P(v_B)-P(x)-P'(x)\Delta|
       \le C B_1B_2 d e^{\theta\max(x,v_B)}\frac{\Delta^2}{x}.
\]
Indeed $x+t\Delta=(1-t)x+t v_B\ge(1-t)x$, so
$(1-t)/(x+t\Delta)\le1/x$ in the remainder integral. This also
handles $v_B=0$ by a limit; there is no singular endpoint error.
Combining with \eqref{biasvariance} and the first derivative estimate
in \eqref{derivatives}, the weighted averaging error is at most
\[
 C B_1B_2e^{\theta\alpha s/(1-\alpha)}
 \left[
  \frac{\sqrt d\,\sqrt s(s+1)}m
  +\frac{d(s+1)}m+\frac{d s(s+1)^2}{m^2}
 \right].
\]
Multiply by $\overline A$, use \eqref{abar} and
\eqref{weightedmoments}, and then multiply by $\ell+1$.
Its contribution is at most $CB_1B_2d/m$.

For rows with $s>m/256$, a coarser estimate suffices. Under a uniformly
chosen subset, $\E_Bv_B=x$ and
\[
 \E_B(v_B-x)^2
 =\frac{\ell}{N(m-1)}\frac1m\sum_{j=1}^m(\ell u_j-x)^2
 \le C(s+1)^2/m.
\]
Since $0\le A_B\le1$, Cauchy--Schwarz gives
\[
 \E_B A_B|v_B-x|
 \le e^{-\alpha s/2} C(s+1)/\sqrt m.
\]
The nonsingular first-derivative bound is
$CB_1B_2d e^{\theta\alpha s/(1-\alpha)}$ on the relevant interval.
As $m/256<s\le m$ and $d\le\sqrt m$, the integrated error, even
including the factor $\ell+1$, is at most $CB_1B_2e^{-cm}$.
Together with \eqref{exactpolarization}, this proves (*).
\end{proof}

\section*{Endpoint tests with an exponential margin}
For $d\ge1$ define
\[
 T_d(y)=\sum_{j=0}^{d-1}(1-j/d)L_j(y),\qquad
 R_d(z)=\left(\frac{2}{\theta d}-z\right)T_d(\theta z)^2.
\]
Laguerre orthogonality and the three-term recurrence give
\[
 \int_0^\infty e^{-y}T_d(y)^2\,dy
   =\frac{(d+1)(2d+1)}{6d},\qquad
 \int_0^\infty ye^{-y}T_d(y)^2\,dy
   =\frac{d+1}{2d}.
\]
Thus the mean of $z$ under $e^{-\theta z}T_d(\theta z)^2\,dz$
is $3/[\theta(2d+1)]\le3/(2\theta d)$.
The density ratio $b_\ell(z)/e^{-\theta z}$ is nonincreasing,
including its drop to zero beyond $\ell$. Multiplication by this
ratio cannot increase the mean: this follows directly by integrating
$(z-z')(r(z)-r(z'))\le0$ against two copies of the original measure.
Consequently
\[
 \beta_\ell(R_d)
 \ge\frac1{2\theta d}\beta_\ell(T_d(\theta z)^2).
\]
For $0\le z\le1/(4d)$ and $j<d$,
$|L_j(\theta z)-1|\le e^{j\theta z}-1<1/3$.
Hence $T_d(\theta z)\ge d/3$ there. Moreover $b_\ell(z)\ge3/4$
on that interval, by $(1-z/\ell)^\ell\ge1-z$.
It follows that
\begin{equation}\label{positiveR}
 \beta_\ell(R_d)\ge\frac1{96\theta}.
\end{equation}

The Laguerre recurrence also gives the useful exact identity
\[
 yT_d(y)=\frac1d\sum_{j=0}^{d-1}L_j(y)-L_d(y).
\]
Consequently $R_d=T_d(\theta z)S_d(z)$, where
\[
 S_d(z)=\frac1{\theta d}\sum_{j=0}^{d-1}(1-2j/d)L_j(\theta z)
                      +\frac1\theta L_d(\theta z).
\]
The sums of absolute Laguerre coefficients are at most $d$ for $T_d$
and $2/\theta$ for $S_d$.
Choose $d=\lfloor\kappa\sqrt m\rfloor$, with a sufficiently small
absolute $0<\kappa\le1$. Lemma~\ref{comparison} bounds the error for
$R_d$ by
\[
 C\left(\frac{d^2}{m}+d e^{-cm}\right).
\]
By first choosing $\kappa$ small and then $m$ large, this is smaller
than half the positive constant in \eqref{positiveR}.
Thus $\E_\mu R_d(x)>0$. Since $R_d$ is nonpositive for
$x\ge2/(\theta d)$, some row of positive $\mu$-mass lies below
that threshold. The last row minimizes $s$ and $x$, so
\begin{equation}\label{lastx}
 x_K<2/(\theta d).
\end{equation}

Put $r=\lfloor d/8\rfloor$ and
\[
 S=\sum_{j=0}^{r}L_j(\theta x_K)^2,
 \qquad Q_*(z)=\frac1S\sum_{j=0}^{r}
              L_j(\theta x_K)L_j(\theta z).
\]
Then $Q_*(x_K)=1$. From \eqref{lastx} we have
$j\theta x_K\le1/4$, so $|L_j(\theta x_K)-1|<1/3$.
In particular $S\ge c d$, and the sum of absolute Laguerre
coefficients of $Q_*$ is bounded by an absolute constant.
Since $b_\ell(z)\le2e^{-z}\le2e^{-\theta z}$,
orthogonality gives
\[
 \beta_\ell(Q_*^2)\le\frac2{\theta S}\le C/d.
\]
Lemma~\ref{comparison}, uniformly also for this data-dependent
polynomial, implies
\[
 \mu(K)\le\E_\mu Q_*(x)^2
 \le \frac Cd+C\left(\frac d m+e^{-cm}\right)
 \le C/d.
\]
Finally $s_K=x_K/\alpha=O(1/d)$. For every $B$,
$A_B(K)=\prod_{j\in B}(1-u_j(K))\ge1-s_K\ge1/2$ for large $m$.
Hence
\[
 (\ell+1)w_K/2\le\mu(K)\le C/d,
 \qquad w_K\le C m^{-3/2}.
\]
Reversing the rows and replacing every $p_j$ by $1-p_j$ gives the
same bound for $w_1$. For a distinguished uniform die, $w_q=1/K$;
the standard comparison tensor has $m=n-1$ and the stated mixed moments.
This proves Theorem~\ref{main}.
\end{document}
